\hspace{3mm}

\centerline{\bf \Huge Налоговые вычеты}

\begin{flushright}
    \scriptsize{ВацIун рикIелай булахар фида,
                \\ гьуьлуьн рикIелай – вацIар.
                \\Лезгинская мудрость.}
\end{flushright}

Некоторые задачи из этого листочка мы разберем на занятии.

\textit{Определение}. 
Ненулевой остаток $a$ называется \textit{квадратичным вычетом} по модулю $p$, если уравнение $x^2 \equiv a \ (mod \ p)$. В иных случаях его называют \textit{квадратичным невычетом.} 

Далее во всех задачах будем считать, что $p$ простое число.

\problem{1} 
Докажите, что квадратичных вычетов и невычетов по модулю $p$ поровну.

\problem{2} 
Докажите, что произведение произведение квадратичных вычетов и невычетов ведёт себя так же как и произведение чётных и нечётных чисел. Говоря очень умными словами, они образуют мультипликативную подгруппу индекса 2 в $\mathbb{Z}_p^*$.

\problem{3}
\textbf{Критерий Эйлера.} Ненулевой остаток $a$ является квадратичным вычетом по модулю $p$ тогда и только тогда, когда:

${\displaystyle a^{\frac {p-1}{2}}\equiv 1{\pmod {p}},}$

\noindentи является квадратичным невычетом по модулю p тогда и только тогда, когда

${\displaystyle a^{\frac {p-1}{2}}\equiv -1{\pmod {p}}.}$

В связи с задачей 2 в математике появилась такая вещь как \textit{символ Лежандра}. Он обозначается через ${\displaystyle \textstyle \left({\frac {a}{p}}\right)}$ и определяется следующим образом:

${\displaystyle \textstyle \left({\frac {a}{p}}\right)}$ = 0, 
если $a$ делится на $p$;

\hspace{3mm}

${\displaystyle \textstyle \left({\frac {a}{p}}\right)}$ = 1,
если $a$ является квадратичным вычетом по модулю $p$.

\hspace{3mm}

${\displaystyle \textstyle \left({\frac {a}{p}}\right)}$ = -1,
если $a$ не является квадратичным вычетом по модулю $p$.

\hspace{3mm}

\problem{4} \textit{Критерий Эйлера при $a = -1$}. 

\noindent
Сформулируйте в терминах остатков на 4, когда -1 является квадратичным вычетом по модулю $p$.

\problem{5}
\textit{Критерий Эйлера при $a = 2$}.

\noindent
Докажите, что при $a=2$ верна формула ${\displaystyle \left({\frac {2}{p}}\right)=(-1)^{(p^{2}-1)/8}.}$

\newpage

\hspace{3mm}

\problem{6}
\textbf{Теорема о корне из -1.}
Число $-1$ является квадратичным вычетом по модулю просто число $p$ тогда и только тогда, когда $p \equiv 1 \ (mod \ 4)$

\problem{7}
\textbf{Теорема Жирара.} 
Пусть $p$ простое число, дающее остаток 3 при делении на 4. Тогда сумма квадратов двух целых чисел делится на $p$ тогда и только тогда, когда они оба делятся на $p$.

\problem{8}
Докажите, что выражение $\dfrac{x^2 + 1}{y^2 - 5}$ никогда не является целым числом при $x, y > 2.$